\begin{center}
    {\LARGE \textbf{2017无机化学}} \\[2ex]
    {\large 时量180分钟 \quad 满分100分}
\end{center}

\vspace{0.5cm}
\noindent\textbf{注意事项:}
\begin{enumerate}[label=\arabic*., parsep=0pt, itemsep=0pt]
    \item 答题前,考生先将自己的姓名、学号、考场号、座位号填写在答题卡上,将条形码准确粘贴在考生信息条形码粘贴区。
    \item 考试结束后,将本试卷与答题纸一并收回。
    \item 回答各个问题时,将答案写在答题纸上,写在本试卷上无效。
\end{enumerate}

\vspace{0.5cm}
\noindent\textbf{一、选择题（20 pts.）}

\begin{enumerate}[label=\arabic*., itemsep=1.5ex]
    \item 下列各对元素中，性质最相似的是
    \begin{tasks}(4)
        \task V和Ta
        \task Hf和Zr
        \task Se和Te
        \task Fe和Ni
    \end{tasks}
    \item 下列氧化还原电对的电极电势值不随$pH$值变化的是
    \begin{tasks}(4)
        \task \ce{MnO4^{-}/Mn^{2+}}
        \task \ce{H2O2/H2O}
        \task \ce{O2/H2O2}
        \task \ce{S2O8^{2-}/SO4^{2-}}
    \end{tasks}
    \item 外围电子构型为$4f^75d^16s^2$的元素在周期表中的位置是
    \begin{tasks}(4)
        \task 第四周期ⅦB族
        \task 第五周期ⅢB族
        \task 第六周期ⅢB族
        \task 第六周期ⅥB族
    \end{tasks}
    \item 在金属晶体的面心立方密堆积结构中，金属原子的配位数为
    \begin{tasks}(4)
        \task 4
        \task 6
        \task 8
        \task 12
    \end{tasks}
    \item 已知某金属离子配合物的磁矩为$4.90\,\text{B.M.}$，而同一氧化态的该金属离子形成的另一配合物，其磁矩为零，则此金属离子可能为
    \begin{tasks}(4)
        \task \ce{Cr(III)}
        \task \ce{Mn(II)}
        \task \ce{Fe(II)}
        \task \ce{Mn(III)}
    \end{tasks}
    \item 下列各组硫化物中，颜色基本相同的是
    \begin{tasks}(4)
        \task \ce{PbS},\ce{Bi2S3},\ce{CuS}
        \task \ce{Ag2S},\ce{HgS},\ce{SnS2}
        \task \ce{CdS},\ce{As2S3},\ce{SnS}
        \task \ce{ZnS},\ce{MnS},\ce{Sb2S3}
    \end{tasks}
    \item 下列羰基化合物中，不符合EAN规则的是
    \begin{tasks}(4)
        \task \ce{[Fe(CO)5]}
        \task \ce{[Co2(CO)8]}
        \task \ce{[Co(CO)4]}
        \task \ce{[Ni(CO)4]}
    \end{tasks}
    \item 下列物质中，可定量吸收\ce{CO}气体的是
    \begin{tasks}(4)
        \task \ce{PdCl2}溶液
        \task 石灰
        \task \ce{KMnO4}溶液
        \task \ce{Fe2O3}
    \end{tasks}
    \item 下列氧化物与浓盐酸反应有氯气放出的是
    \begin{tasks}(4)
        \task \ce{Fe2O3}
        \task \ce{V2O5}
        \task \ce{TiO2}
        \task \ce{SnO2}
    \end{tasks}
    \item 下列各组元素中，原子的第一电离能递增顺序正确的是
    \begin{tasks}(4)
        \task \ce{Na} $<$ \ce{Mg} $<$ \ce{Al}
        \task \ce{He} $<$ \ce{Ne} $<$ \ce{Ar}
        \task \ce{Si} $<$ \ce{P} $<$ \ce{As}
        \task \ce{B} $<$ \ce{C} $<$ \ce{N}
    \end{tasks}
\end{enumerate}

\vspace{0.5cm}
\noindent\textbf{二、解下列实验现象，写相关的反应方程式（30 pts.）}

\begin{enumerate}[label=\arabic*., start=1, itemsep=1.5ex]
    \item 分别向\ce{FeCl2}和\ce{NaHCO3}溶液中加入碘水，碘水均不褪色。但向二者的混合溶液中加入碘水，则碘水褪色。
    \item 将\ce{As2O3}与金属锌、盐酸混合有气体产生，将该气体通过一玻璃管（尾端有吸收装置），并在入口处加热，则玻璃管壁出现有金属光泽的黑色物质。
    \item 测得某硼酸溶液的$pH$值为6，加入一些甘油（\ce{C3H5(OH)3}）后，测得溶液的$pH$值变为5。
    \item 向\ce{Hg2(NO3)2}溶液中不断滴加\ce{KI}溶液，先有黄色沉淀生成，继而沉淀颜色加深，最后得到黑色沉淀和无色溶液。
    \item 少量\ce{Na2S2O3}溶液和\ce{AgNO3}溶液反应生成白色沉淀，沉淀逐渐变黄变棕色，最后变成黑色。白色沉淀可溶于过量\ce{Na2S2O3}溶液，而黑色沉淀却不溶于过量\ce{Na2S2O3}溶液。
\end{enumerate}

\vspace{0.5cm}
\noindent\textbf{三、分离、鉴别与制备（20 pts.）}

\begin{enumerate}[label=\arabic*., start=1, itemsep=1.5ex]
    \item 用化学反应方程式表示工业上由浓缩的海水中提取单质溴的过程。
    \item 4个试剂瓶中分别装有\ce{Na4P2O7}、\ce{NaPO3}、\ce{Na2HPO4}和\ce{NaH2PO2}固体，试加以鉴别。
    \item 如何在不引入杂质的情况下，将粗\ce{ZnSO4}溶液中含有的\ce{Fe^2+}，\ce{Fe^{3+}}和\ce{Cu^{2+}}杂质除去？
    \item 设计方案分离混合溶液中的离子。\ce{Cr^{3+}}、\ce{Al^{3+}}、\ce{Mn^{2+}}、\ce{Ag^+}
\end{enumerate}

\vspace{0.5cm}
\noindent\textbf{四、回答下列各题（65 pts.）}

\begin{enumerate}[label=\arabic*., start=1, itemsep=1.5ex]
    \item 写出下列物质的化学式。
    \begin{enumerate}[label=(\arabic*), itemsep=1ex]
        \item 萤石
        \item 铅丹
        \item 金红石
        \item 海波
        \item 芒硝
        \item 乙硼烷
        \item 石英
    \end{enumerate}
    \item 试说明为什么氟的电子亲和能比氯低，但\ce{F2}却比\ce{Cl2}活泼。
    \item \ce{[Ni(CO)4]}和\ce{[Ni(CN)4]^{2-}}都是抗磁性，两者在结构上是否相同？试用配位化合物的价键理论给予说明。
    \item 比较大小，并简要说明理由。
    \begin{enumerate}[label=(\arabic*), itemsep=1ex]
        \item 氧化性：\ce{HClO4}，\ce{H2SO4}，\ce{H3PO4}、\ce{H4SiO4}
        \item 热稳定性：\ce{Na2CO3}，\ce{MgCO3}，\ce{K2CO3}，\ce{PbCO3}
    \end{enumerate}
    \item 简要回答。
    \begin{enumerate}[label=(\arabic*), itemsep=1ex]
        \item 为什么不能采取高温浓缩的办法制得\ce{NaHSO3}晶体？
        \item 为什么\ce{CCl4}遇水不发生水解，而\ce{BCl3}和\ce{SiCl4}却强烈水解？
    \end{enumerate}
    \item 判断下列分子或离子的空间构型，并说明中心原子的杂化方式。
    \begin{enumerate}[label=(\arabic*), itemsep=1ex]
        \item \ce{IF3}
        \item \ce{SO3^{2-}}
        \item \ce{XeOF4}
        \item \ce{[Fe(CO)5]}
    \end{enumerate}
    \item 测定水中溶解氧量的方法是：在隔绝空气的条件下，向一定量的水样中加入适量的\ce{MnCl2}及过量的\ce{NaOH}溶液；然后用过量的\ce{H2SO4}酸化，再加入过量的\ce{KI}溶液。充分反应后用标准的\ce{Na2S2O3}溶液滴定生成的\ce{I2}。试写出有关的反应方程式，并给出\ce{O2}和消耗\ce{Na2S2O3}的物质的量之比。
    已知相关电对的电极电势值为：$E^\ominus(\ce{MnO(OH)2/Mn(OH)2}) = -0.05\mathrm{V}$；$E^\ominus(\ce{MnO2/Mn^{2+}}) = 1.23\mathrm{V}$；$E^\ominus(\ce{O2/OH^-}) = 0.401\mathrm{V}$；$E^\ominus(\ce{I2/I^-}) = 0.536\mathrm{V}$；$E^\ominus(\ce{S4O6^{2-}/S2O3^{2-}}) = 0.08\mathrm{V}$。
    \item 已知：
    \[
    \begin{array}{ll}
    \ce{Cu^{2+} + 2e^- = Cu} & E^\ominus = 0.34\mathrm{V} \\
    \ce{Cu^+ + e^- = Cu} & E^\ominus = 0.52\mathrm{V} \\
    \ce{Cu^{2+} + Br^- + e^- = CuBr} & E^\ominus = 0.64\mathrm{V}
    \end{array}
    \]
    试求\ce{CuBr}的$K_{sp}$。
\end{enumerate}

\vspace{0.5cm}
\noindent\textbf{五、推断题（15 pts.）}

\begin{enumerate}[label=\arabic*., start=1, itemsep=1.5ex]
    \item 化合物(A)为棕褐色不溶于水的固体，将(A)溶解在浓盐酸中，得到溶液(B)和黄绿色气体(C)，(C)通过热氢氧化钾溶液，被吸收后生成溶液(D)，(D)用酸中和并以\ce{AgNO3}溶液处理生成白色沉淀(E)，(E)不溶于硝酸，但可溶于氨水。若在(B)溶液中加入氢氧化钠溶液，得到粉红色沉淀(F)，将(F)用\ce{H2O2}处理，则又得到(A)。若将(B)用\ce{KNO2}的醋酸溶液处理，则生成黄色沉淀(G)和气体(H)，将(H)通过\ce{FeSO4}溶液，则溶液变为暗棕色，但没有沉淀生成。试给出(A)、(B)、(C)、(D)、(E)、(F)、(G)和(H)所代表的物质的化学式，并用化学反应方程式表示各过程。
\end{enumerate}